\section{Introdução}
\label{sec:introducao}

\paragraph{}A simulação numérica consolidou-se nas últimas décadas como o terceiro pilar da ciência contemporânea, posicionando-se ao lado da teoria e da experimentação. Na Engenharia Mecânica, especificamente, a Dinâmica dos Fluidos Computacional (CFD - \textit{Computational Fluid Dynamics}) e a análise de transferência de calor por métodos numéricos permitem o estudo de fenômenos complexos que seriam proibitivos, perigosos ou impossíveis de serem replicados em bancada experimental \cite{Maliska2004, Versteeg2007}. Desde o projeto de turbomáquinas até a análise de dispersão de poluentes, a capacidade de prever o comportamento de fluidos através da resolução das equações de Navier-Stokes transformou a indústria e a academia.

\paragraph{}No entanto, o acesso a essas ferramentas poderosas historicamente esbarra em barreiras significativas. Softwares comerciais de referência, como ANSYS Fluent, COMSOL Multiphysics ou SimScale, embora robustos, impõem custos de licenciamento elevados e exigem hardware dedicado de alta performance (HPC) para execução eficiente. No contexto educacional e de pesquisa inicial, essas barreiras dificultam a democratização do conhecimento em métodos numéricos. Estudantes e engenheiros frequentemente precisam lidar com curvas de aprendizado íngremes e instalações complexas de ambiente antes de rodar sua primeira simulação de escoamento \cite{Anjos2013}.

\paragraph{}Este trabalho propõe uma mudança de paradigma ao trazer a Mecânica Computacional para a web. Aproveitando o avanço recente das tecnologias de navegador, especificamente o \textit{WebAssembly} e o projeto \textit{PyScript} \cite{PyScript2023}, torna-se possível executar códigos científicos em Python — a \textit{lingua franca} da ciência de dados atual — diretamente no navegador do usuário, sem necessidade de servidores backend ou instalações locais.

\section{Objetivos}

\paragraph{}O objetivo geral deste trabalho é desenvolver e validar o \textbf{MinervaMesh}, uma plataforma computacional baseada na web para a solução de equações diferenciais parciais governantes de problemas de Engenharia Mecânica, operando inteiramente no lado do cliente (\textit{client-side}).

\paragraph{}Os objetivos específicos são:
\begin{enumerate}
    \item Implementar métodos numéricos clássicos, especificamente o Método dos Elementos Finitos (MEF) e o Método das Diferenças Finitas (MDF), utilizando bibliotecas científicas Python (NumPy, SciPy) rodando via WebAssembly \cite{Harris2020}.
    \item Desenvolver uma interface gráfica interativa que permita a pré-processamento (configuração de malha e parâmetros), solução e pós-processamento (visualização de campos de velocidade, pressão e temperatura) em tempo real.
    \item \textbf{Validar a plataforma com foco principal em Mecânica dos Fluidos}, comparando os resultados numéricos com soluções analíticas e benchmarks estabelecidos para:
    \begin{itemize}
        \item Escoamentos potenciais;
        \item Problemas de advecção-difusão (com estabilização SUPG) \cite{Brooks1982};
        \item Equações de Navier-Stokes para escoamentos laminares incompressíveis.
    \end{itemize}
\end{enumerate}

\section{Justificativa}

\paragraph{}A motivação para este desenvolvimento reside na necessidade de ferramentas ágeis para o ensino e prototipagem rápida em engenharia. A capacidade de rodar uma simulação de escoamento em qualquer dispositivo (incluindo tablets e laptops modestos), apenas acessando uma URL, remove o atrito inicial do aprendizado de CFD. Além disso, ao utilizar tecnologias de código aberto e documentar a implementação dos métodos numéricos \cite{Fish2007, Reddy2005}, o \textit{MinervaMesh} serve não apenas como ferramenta, mas como objeto de estudo para estudantes interessados na implementação de seus próprios códigos de simulação.

\section{Organização do Trabalho}

\paragraph{}Este trabalho está estruturado da seguinte forma:
\begin{itemize}
    \item O \textbf{Capítulo 2} (Fundamentação Teórica) revisa as equações de conservação (Massa, Momento e Energia) \cite{Fox2010, White2006} e apresenta a formulação matemática dos métodos numéricos empregados (MEF e MDF), com ênfase na discretização espacial e temporal.
    \item O \textbf{Capítulo 3} (Desenvolvimento e Arquitetura) detalha a construção da plataforma \textit{MinervaMesh}, discutindo decisões de arquitetura de software, o uso do PyScript e a integração das bibliotecas numéricas no ambiente web.
    \item O \textbf{Capítulo 4} (Resultados e Validação) é o núcleo da verificação do trabalho. Nele, são apresentados os casos de teste, com foco especial nos problemas de \textbf{dinâmica dos fluidos}, demonstrando a precisão através de análises de erro e convergência de malha.
    \item O \textbf{Capítulo 5} (Conclusões) resume as contribuições, limitações encontradas e propõe trabalhos futuros.
\end{itemize}
